% Generated draft: requires mathematical and source review.
\documentclass[11pt]{article}
\usepackage[margin=25mm]{geometry}
\usepackage{fontspec}
\setmainfont{Latin Modern Roman}
\usepackage{amsmath,amssymb,mathtools,unicode-math}
\setmathfont{Latin Modern Math}
\usepackage{enumitem,needspace,xurl,hyperref}
\setlength{\parindent}{0pt}
\setlength{\parskip}{5pt}
\setlength{\emergencystretch}{3em}
\newcommand{\N}{\mathbb{N}}
\newcommand{\Z}{\mathbb{Z}}
\newcommand{\Q}{\mathbb{Q}}
\newcommand{\R}{\mathbb{R}}
\newcommand{\C}{\mathbb{C}}
\newcommand{\F}{\mathbb{F}}
\newcommand{\A}{\mathbb{A}}
\newcommand{\PP}{\mathbb{P}}
\newcommand{\E}{\mathbb{E}}
\newcommand{\eps}{\varepsilon}
\newcommand{\dd}{\,\mathrm{d}}
\begin{document}
\section*{UnsolvedMath ID: 20002357}
\textbf{Source problem code:} catalogue entry
\subsection{Definitions and mathematical statement}
Let \(E/\Q\) be an elliptic curve and \(G=E(\Q)\). Its Mordell--Weil rank is \(\dim_\Q(G\otimes_\Z\Q)\); assume this rank is two, so \(G\cong\Z^2\oplus F\) for a finite abelian torsion group \(F\). A subgroup \(H\le G\) has rank one if \(\dim_\Q(H\otimes_\Z\Q)=1\). Interpreting definability in the field \(\Q\) with its ring language, rational parameters, and an algebraic coordinate coding of \(E(\Q)\), Shlapentokh's Question 20 asks
\[
 \forall E/\Q\text{ of rank }2\quad
 \exists H\le E(\Q)\text{ of rank }1
 \quad H\text{ is existentially definable over }\Q.
\]
Over \(\Q\), inequations can be converted to existential equations by inverses; finite coordinate charts and the identity point must be included when coding the projective curve. The short source does not itself specify a language or parameter convention; the ring-language interpretation is the one used by the aid.

Fix an infinite-order \(P\in G\) and write \(H_P=\langle P\rangle\). For each fixed positive integer \(m\), define the congruence tube
\[
 T_m=H_P+[m]G
 =\bigcup_{r=0}^{m-1}\bigl([r]P+[m]G\bigr).
\]
Thus membership in \(T_m\) is expressed by the finite disjunction
\[
 X\in G\quad\wedge\quad
 \bigvee_{r=0}^{m-1}\exists Q\in G\ (X=[r]P+[m]Q).
\]
Using the algebraic group law, each fixed tube is positive-existentially definable over \(\Q\), with \(P\) fixed. This means one finite formula for each fixed \(m\), not an asserted single Diophantine formula with a variable integer multiplication index.

The supplied aid asserts the exact limit description
\[
 H_P=\bigcap_{N\ge1}\bigl(H_P+[N!]G\bigr),\qquad
 H_P+[(N+1)!]G\subseteq H_P+[N!]G.
\]
The finite-generation hypothesis matters here: \(G/H_P\) is finitely generated, and its intersection of factorial multiples is zero, including its finite torsion part. This infinite conjunction is not a first-order existential definition. Moreover, if a nonempty finite expression is built from tubes \(T_{m_1},\ldots,T_{m_s}\) using only intersections and unions, it contains \([L]G\), where \(L=\operatorname{lcm}(m_1,\ldots,m_s)\). Hence it contains a rank-two subgroup and cannot itself be a rank-one subgroup. This obstruction concerns those finite positive combinations, not arbitrary field-language formulas.

The independent pure-group obstruction reported by the aid is: for \(G\cong\Z^2\oplus F\), no subgroup of rank one is definable in one free group variable by a positive-existential formula in \(\{0,+,-\}\), even with parameters from \(G\). Such formulas use existential quantifiers, conjunctions, and disjunctions of group equations, not arbitrary negations or field arithmetic. The assertion does not cover the intended ring-language problem, whose quantified witnesses can range over rational field coordinates. Thus the aid's negative pure-group result and factorial-limit description leave the original existential field-arithmetic question open.
\subsection{Short English statement}
Does every rank-two rational elliptic curve have an existentially definable rank-one subgroup? The aid describes a cyclic subgroup as an infinite intersection of individually Diophantine congruence tubes, while ruling out finite positive tube combinations and positive-existential definitions using only the abstract group language.
\subsection{Sources}
\begin{itemize}
\item Local registry, permanent UnsolvedMath ID \(20002357\), exact public code \texttt{catalogue entry}; canonical location \texttt{aim-logic-notes.json}, \texttt{notes[132]}.
\item AIM, \emph{Extensions of Hilbert's tenth problem}, Question 20 attributed to Shlapentokh; supplied reference \url{https://aimath.org/WWN/hilberts10th/hilberts10th.pdf}. The verification field reports an undamaged official statement; no inspection was performed here. No proposer year, original remarks, or literature citation is supplied.
\item Supplied research attempt 1: \texttt{partial\_result}, run status \texttt{open}, rigor label \texttt{proved}, low novelty confidence; registry status \texttt{partially\_solved}, classification \texttt{PARTIAL-PROGRESS}, created and updated 14 August 2026. The tube and pure-group assertions have no additional primary bibliographic citation in the record.
\item The aid explicitly states that it is machine generated, not peer reviewed, and requires independent verification before citation.
\end{itemize}
\end{document}
